\section{Background}
\label{sec:background}

\subsection{MC-PILCO}
\label{sec:background:mcpilco}

\mcpilco{}~\cite{mcpilco} is a model-based RL framework for continuous control that
couples a GP dynamics model with a Monte Carlo (MC) policy-gradient optimiser.
The GP is fit to observed transitions $(\mathbf{x}_t, \mathbf{u}_t) \to \mathbf{x}_{t+1}$
and maintains a posterior over residual corrections to a known nominal model.
At each policy update step, $M$ particles are propagated forward through the GP posterior
via the reparameterisation trick: each particle draws a function sample from the GP and
integrates the resulting ODE, yielding a set of stochastic rollouts that can be
differentiated with respect to policy parameters $\theta$.
The policy gradient is then

\begin{equation}
  \nabla_\theta J(\theta) = \nabla_\theta \mathbb{E}\!\left[
    \sum_{t=0}^{T} c(\tilde{\mathbf{x}}_t)
  \right] \approx \frac{1}{M}\sum_{m=1}^{M} \nabla_\theta \sum_{t=0}^{T}
  c(\tilde{\mathbf{x}}_t^{(m)}),
  \label{eq:mcpilco_gradient}
\end{equation}

where the reparameterisation ensures that the Monte Carlo samples are differentiable
through the GP posterior.
Sparse approximation via Subset of Data (SOD) keeps inference tractable as the dataset
grows across trials.

\subsection{MC-PILOT Modifications}
\label{sec:background:mcpilot}

\mcpilot{}~\cite{turcato2025mcpilot} adapts \mcpilco{} for the single-shot throwing
problem with two targeted modifications.

\paragraph{Augmented state.}
The robot state $\mathbf{x} = [\mathbf{p},\, \dot{\mathbf{p}}] \in \mathbb{R}^6$
(position and velocity) is augmented with the target location
$\mathbf{P} = (P_x, P_y) \in \mathbb{R}^2$ to form
$\tilde{\mathbf{x}} = [\mathbf{x},\, \mathbf{P}] \in \mathbb{R}^8$.
The GP learns per-timestep velocity corrections
$[\Delta v_x, \Delta v_y, \Delta v_z]$ conditioned on $\tilde{\mathbf{x}}$, enabling
geometry-conditioned residual compensation across the target workspace.

\paragraph{Single-shot policy.}
Because the projectile is in free flight after release, feedback control is neither
possible nor necessary.
The policy collapses to a single map from target position to release speed:

\begin{equation}
  \pi_\theta(\mathbf{P}) = \frac{u_M}{2}\!\left(
    \tanh\!\left(
      \sum_{i=1}^{N_b} w_i \exp\!\left(
        -\frac{\|\mathbf{a}_i - \mathbf{P}\|^2}{\ell_s^2}
      \right)
    \right) + 1
  \right),
  \label{eq:policy}
\end{equation}

where $u_M$ is the maximum allowable speed, $\{(\mathbf{a}_i, w_i)\}_{i=1}^{N_b}$ are
the $N_b$ RBF centres and weights (the trainable parameters $\theta$), and $\ell_s$ is
the lengthscale that controls the spatial extent of each basis function.
The $\tanh(\cdot) + 1$ construction squashes the output to $(0,\, u_M)$, enforcing
physical speed limits.

\paragraph{Landing cost.}
Performance is measured at the final timestep $T$ (the moment of landing) by a
Gaussian-shaped cost centred on the target:

\begin{equation}
  c(\tilde{\mathbf{x}}_T) = 1 - \exp\!\left(
    -\|\mathbf{p}_T - \mathbf{P}\|^2 \cdot \Sigma_c^{-1}
  \right),
  \label{eq:cost}
\end{equation}

where $\mathbf{p}_T \in \mathbb{R}^2$ is the ball landing position and $\Sigma_c$
controls the width of the acceptance region.
The policy objective is $J(\theta) = \mathbb{E}[c(\tilde{\mathbf{x}}_T)]$, minimised
by gradient descent through the MC rollouts.

\subsection{Aerodynamic Drag Model}
\label{sec:background:drag}

Ball trajectories in \mcpilot{} include aerodynamic drag.
Following Eq.~35 of~\cite{turcato2025mcpilot}, the drag force is

\begin{equation}
  \mathbf{F}_D = -\tfrac{1}{2}\rho\, C_D(\mathrm{Re})\, A\,
    \|\dot{\mathbf{p}}\|\, \dot{\mathbf{p}},
  \label{eq:drag}
\end{equation}

where $\rho$ is air density, $A = \pi r^2$ is the ball cross-section, and the drag
coefficient $C_D(\mathrm{Re})$ depends on the Reynolds number
$\mathrm{Re} = \|\dot{\mathbf{p}}\| \cdot 2r / \nu$, with $r$ the ball radius and $\nu$
the kinematic viscosity of air.
For the tennis ball used in our experiments ($m = 0.0577$\,kg, $r = 0.0327$\,m) and
typical release speeds below 4\,m/s, $\mathrm{Re}$ remains in the subcritical regime
where $C_D \approx 0.47$, so we use a constant coefficient.

\subsection{Omitted Components}
\label{sec:background:omissions}

We intentionally exclude two components from our reproduction.
First, the gripper delay estimation module, which in the original paper uses Bayesian
Optimisation to calibrate a timing offset between the commanded and actual release
instant~\cite{turcato2025mcpilot}: we instead model the effect of release timing
uncertainty directly as a stochastic perturbation in our noise study
(Section~\ref{sec:study2}), making the BO calibration step unnecessary.
Second, all hardware-specific components (ROS communication, Gazebo physics server,
Franka Panda driver) are replaced by a self-contained numpy ballistic integrator
described in Section~\ref{sec:system}.
These omissions are intentional and do not affect the core learning algorithm under
evaluation.
